% GENERATED FILE. Edit canonical lessons, then run npm run generate:books.
% canonical-source-hash: fnv1a64:30ddeafa0001ca75
% canonical-lessons: TE-C488-toku, TE-C488-nilva, TE-C488-giddangi, TE-C488-korata, TE-C488-velam

\chapter{Wholesale and Stock}
\label{ch:te-wholesale-and-stock}
\hlchaptermodality{488}

\begin{chapteropening}
\textbf{By the end of this chapter:} \emph{I can say wholesale, a stock, a storehouse, a shortage and an auction.}

Complete the last lesson of chapter 488: I can say wholesale, a stock, a storehouse, a shortage and an auction.
\end{chapteropening}

\section[\texorpdfstring{\te{టోకు} (ṭōku)}{టోకు (ṭōku)}]{\te{టోకు} (ṭōku) — wholesale}

\label{lesson:TE-C488-toku}



\begin{quote}
Before the new one: say the Telugu for to be available, then the Telugu for to sell out.
\end{quote}

\subsection*{You'll want to know: \te{టోకు}}
\textbf{\te{టోకు}} — \emph{ṭōku} — \textquotedblleft{}wholesale\textquotedblright{}.

\begin{quote}
\te{అతను} \te{టోకు} \te{ధరకు} \te{బియ్యం} \te{కొంటాడు}.
\end{quote}

\begin{quote}
\emph{atanu ṭōku dharaku biyyaṁ koṇṭāḍu.}
\end{quote}

\begin{quote}
He buys rice at the wholesale price.
\end{quote}

\te{టోకు} \te{వ్యాపారి}, a wholesale dealer; its opposite is \te{చిల్లర}, retail.

\begin{cousinweb}
\textbf{\te{టోకు}} (\emph{ṭōku}) is a market word for wholesale whose origin is not settled.
\end{cousinweb}

\begin{grammarlens}[title={What you've built}]
One more word for getting around.
\end{grammarlens}

\subsection*{Your turn}
\begin{itemize}
\raggedright
  \item \emph{Say it:} \emph{ṭōku}
  \item \emph{Say it:} \emph{ṭōku}, once more
  \item \emph{Say it:} \emph{atanu ṭōku dharaku biyyaṁ koṇṭāḍu.} — He buys rice at the wholesale price.
  \item \emph{Recall:} say the Telugu for to be available, then the Telugu for to sell out, then say \emph{ṭōku} again
\end{itemize}

\begin{tcolorbox}[breakable,colback=teal!4,colframe=teal!35!black,title={Before you move on}]
What does \te{టోకు} mean? (\textquotedblleft{}Wholesale\textquotedblright{}.) Say it once more.
\end{tcolorbox}
\section[\texorpdfstring{\te{నిల్వ} (nilva)}{నిల్వ (nilva)}]{\te{నిల్వ} (nilva) — a stock; a store}

\label{lesson:TE-C488-nilva}



\begin{quote}
Before the new one: say the Telugu for to sell out, then the Telugu for wholesale.
\end{quote}

\subsection*{You'll want to know: \te{నిల్వ}}
\textbf{\te{నిల్వ}} — \emph{nilva} — \textquotedblleft{}a stock; a store\textquotedblright{}.

\begin{quote}
\te{దుకాణంలో} \te{పంచదార} \te{నిల్వ} \te{అయిపోయింది}.
\end{quote}

\begin{quote}
\emph{dukāṇaṁlō pañcadāra nilva ayipōyindi.}
\end{quote}

\begin{quote}
The shop has run out of its stock of sugar.
\end{quote}

\te{నిల్వ} \te{చేయు}, to store. \te{నిల్వ} \te{పచ్చడి} is a pickle made to keep.

\begin{cousinweb}
\textbf{\te{నిల్వ}} (\emph{nilva}) is native Dravidian, of the family of \te{నిలుచు}, \textquotedblleft{}to stand, to remain, to be preserved\textquotedblright{} (DEDR 3675): what remains is the stock.
\end{cousinweb}

\begin{grammarlens}[title={What you've built}]
One more word for getting around.
\end{grammarlens}

\subsection*{Your turn}
\begin{itemize}
\raggedright
  \item \emph{Say it:} \emph{nilva}
  \item \emph{Say it:} \emph{nilva}, once more
  \item \emph{Say it:} \emph{dukāṇaṁlō pañcadāra nilva ayipōyindi.} — The shop has run out of its stock of sugar.
  \item \emph{Recall:} say the Telugu for to sell out, then the Telugu for wholesale, then say \emph{nilva} again
\end{itemize}

\begin{tcolorbox}[breakable,colback=teal!4,colframe=teal!35!black,title={Before you move on}]
What does \te{నిల్వ} mean? (\textquotedblleft{}A stock; a store\textquotedblright{}.) Say it once more.
\end{tcolorbox}
\section[\texorpdfstring{\te{గిడ్డంగి} (giḍḍaṅgi)}{గిడ్డంగి (giḍḍaṅgi)}]{\te{గిడ్డంగి} (giḍḍaṅgi) — a storehouse; a godown}

\label{lesson:TE-C488-giddangi}



\begin{quote}
Before the new one: say the Telugu for wholesale, then the Telugu for a stock.
\end{quote}

\subsection*{You'll want to know: \te{గిడ్డంగి}}
\textbf{\te{గిడ్డంగి}} — \emph{giḍḍaṅgi} — \textquotedblleft{}a storehouse; a godown\textquotedblright{}.

\begin{quote}
\te{రైతులు} \te{ధాన్యాన్ని} \te{గిడ్డంగిలో} \te{దాచారు}.
\end{quote}

\begin{quote}
\emph{raitulu dhānyānni giḍḍaṅgilō dācāru.}
\end{quote}

\begin{quote}
The farmers stored the grain in the warehouse.
\end{quote}

A government warehouse for grain is a \te{గిడ్డంగి}.

\begin{cousinweb}
\textbf{\te{గిడ్డంగి}} (\emph{giḍḍaṅgi}) is native Dravidian: DEDR 1525 lists Telugu \emph{giḍḍaṅgi}, \textquotedblleft{}a storehouse, a godown\textquotedblright{}, with Tamil \emph{kiṭaṅku}, \textquotedblleft{}a storehouse\textquotedblright{}.
\end{cousinweb}

\begin{grammarlens}[title={What you've built}]
One more word for getting around.
\end{grammarlens}

\subsection*{Your turn}
\begin{itemize}
\raggedright
  \item \emph{Say it:} \emph{giḍḍaṅgi}
  \item \emph{Say it:} \emph{giḍḍaṅgi}, once more
  \item \emph{Say it:} \emph{raitulu dhānyānni giḍḍaṅgilō dācāru.} — The farmers stored the grain in the warehouse.
  \item \emph{Recall:} say the Telugu for wholesale, then the Telugu for a stock, then say \emph{giḍḍaṅgi} again
\end{itemize}

\begin{tcolorbox}[breakable,colback=teal!4,colframe=teal!35!black,title={Before you move on}]
What does \te{గిడ్డంగి} mean? (\textquotedblleft{}A storehouse; a godown\textquotedblright{}.) Say it once more.
\end{tcolorbox}
\section[\texorpdfstring{\te{కొరత} (korata)}{కొరత (korata)}]{\te{కొరత} (korata) — a shortage}

\label{lesson:TE-C488-korata}



\begin{quote}
Before the new one: say the Telugu for a stock, then the Telugu for a storehouse.
\end{quote}

\subsection*{You'll want to know: \te{కొరత}}
\textbf{\te{కొరత}} — \emph{korata} — \textquotedblleft{}a shortage\textquotedblright{}.

\begin{quote}
\te{వేసవిలో} \te{నీటి} \te{కొరత} \te{ఉంది}.
\end{quote}

\begin{quote}
\emph{vēsavilō nīṭi korata undi.}
\end{quote}

\begin{quote}
In summer there is a shortage of water.
\end{quote}

\te{నీటి} \te{కొరత}, a water shortage; \te{డబ్బు} \te{కొరత}, a lack of money.

\begin{cousinweb}
\textbf{\te{కొరత}} (\emph{korata}) is native Dravidian: DEDR 1851 lists Telugu \emph{koṟata}, \textquotedblleft{}deficiency, want\textquotedblright{}, in the family of \emph{kuṟu}, \textquotedblleft{}short, small\textquotedblright{}.
\end{cousinweb}

\begin{grammarlens}[title={What you've built}]
One more word for getting around.
\end{grammarlens}

\subsection*{Your turn}
\begin{itemize}
\raggedright
  \item \emph{Say it:} \emph{korata}
  \item \emph{Say it:} \emph{korata}, once more
  \item \emph{Say it:} \emph{vēsavilō nīṭi korata undi.} — In summer there is a shortage of water.
  \item \emph{Recall:} say the Telugu for a stock, then the Telugu for a storehouse, then say \emph{korata} again
\end{itemize}

\begin{tcolorbox}[breakable,colback=teal!4,colframe=teal!35!black,title={Before you move on}]
What does \te{కొరత} mean? (\textquotedblleft{}A shortage\textquotedblright{}.) Say it once more.
\end{tcolorbox}
\section[\texorpdfstring{\te{వేలం} (vēlaṁ)}{వేలం (vēlaṁ)}]{\te{వేలం} (vēlaṁ) — an auction}

\label{lesson:TE-C488-velam}



\begin{quote}
Before the new one: say the Telugu for a storehouse, then the Telugu for a shortage.
\end{quote}

\subsection*{You'll want to know: \te{వేలం}}
\textbf{\te{వేలం}} — \emph{vēlaṁ} — \textquotedblleft{}an auction\textquotedblright{}.

\begin{quote}
\te{పాత} \te{బస్సులను} \te{వేలంలో} \te{అమ్మారు}.
\end{quote}

\begin{quote}
\emph{pāta bassulanu vēlaṁlō ammāru.}
\end{quote}

\begin{quote}
The old buses were sold at auction.
\end{quote}

\te{వేలం} \te{వేయు}, to auction; \te{వేలంపాట}, the bidding.

\begin{cousinweb}
\textbf{\te{వేలం}} (\emph{vēlaṁ}) is usually traced to Portuguese \emph{leilão}, \textquotedblleft{}an auction\textquotedblright{}, from the trade of the Portuguese on the Indian coasts.
\end{cousinweb}

\begin{grammarlens}[title={What you've built}]
One more word for getting around.
\end{grammarlens}

\subsection*{Your turn}
\begin{itemize}
\raggedright
  \item \emph{Say it:} \emph{vēlaṁ}
  \item \emph{Say it:} \emph{vēlaṁ}, once more
  \item \emph{Say it:} \emph{pāta bassulanu vēlaṁlō ammāru.} — The old buses were sold at auction.
  \item \emph{Recall:} say the Telugu for a storehouse, then the Telugu for a shortage, then say \emph{vēlaṁ} again
\end{itemize}

\begin{tcolorbox}[breakable,colback=teal!4,colframe=teal!35!black,title={Before you move on}]
What does \te{వేలం} mean? (\textquotedblleft{}An auction\textquotedblright{}.) Say it once more.
\end{tcolorbox}
